% Table T03 -- Recording Duration and Sampling Summary
% Experiment: n/a (not experiment-bound)
% Objective: O1 (corpus definition)
% Source: outputs/01_dataset_audit/recording_duration_summary.csv; outputs/01_dataset_audit/sampling_rate_summary.csv
% Generated by PV-MEPCG / PulseVision at 2026-08-30T10:15:33.877856+00:00
\begin{table}[htbp]
  \centering
  \caption{Recording length in seconds over the supervised records of each corpus, and the sampling rate conversion applied to it. All four corpora are resampled to a common 2 kHz working rate; the durations are of the original recordings and are unchanged by resampling. Total corpus hours are in T01 and cover all files, not only the supervised subset summarised here.}
  \label{tab:t03}
  \begin{tabular}{lrrrrrrrrrrrr}
    \toprule
    ID & Dataset & Records & Min (s) & P25 (s) & Median (s) & P75 (s) & Max (s) & Mean (s) & SD (s) & Native fs (Hz) & Working fs (Hz) & Conversion \\
    \midrule
    D1 & PhysioNet 2016 & 3,240 & 5.31 & 12.90 & 20.83 & 30.58 & 122.00 & 22.46 & 12.39 & 2,000 & 2,000 & none \\
    D2 & PASCAL set\_a & 124 & 0.94 & 7.94 & 8.88 & 9.00 & 9.00 & 8.01 & 1.64 & 44,100 & 2,000 & decimation \\
    D3 & PASCAL set\_b & 461 & 0.76 & 2.97 & 5.09 & 9.31 & 27.87 & 6.48 & 4.48 & 4,000 & 2,000 & decimation \\
    D4 & CirCor 2022 & 3,163 & 5.15 & 19.06 & 21.46 & 29.39 & 64.51 & 22.87 & 7.28 & 4,000 & 2,000 & decimation \\
    \bottomrule
  \end{tabular}
  \par\footnotesize The duration extremes are real and both are handled by the extractor: PASCAL set\_b reaches 0.76 s, too short for a five-level DWT, and PhysioNet reaches 122.0 s.
\end{table}
